\documentclass[10pt,a4paper]{amsart}
\usepackage[margin=19mm]{geometry}
\usepackage{amsmath,amssymb,mathtools}
\usepackage[scr=rsfs,cal=euler]{mathalfa}
\usepackage{xfrac}
\usepackage[unicode,hidelinks,bookmarks=false]{hyperref}
\newcommand{\ie}{i.e.\ }
\newcommand{\cf}{cf.\ }
\newcommand{\resp}{respectively\ }
\newcommand{\Subsection}{\subsection}
\providecommand{\nobreakdash}{\nobreak\hbox{-}}
\setlength{\emergencystretch}{2em}
\multlinegap0pt
\usepackage{mathtools}
\usepackage{amssymb}
\usepackage{xcolor}
\usepackage{tikz}
\usepackage{algorithm}\usepackage{algpseudocode}
\newtheorem{problem}{Problem}
\newtheorem{theorem}{Theorem}
\newtheorem{proposition}{Proposition}
\usepackage{cleveref}
\crefname{problem}{Problem}{Problems}
\crefname{theorem}{Theorem}{Theorems}
\crefname{proposition}{Proposition}{Propositions}
\newcommand{\Mset}{\langle M\rangle}
\title{\LARGE \bf Iterative graph lifting for automatic design of path-complete stability certificates}
\author{L\'ea Ninite, Rapha\"el M. Jungers}
\date{Source: arXiv:2607.00637v2 (CC BY 4.0)}
\begin{document}
\maketitle

\noindent\emph{Source and licence.} \LARGE \bf Iterative graph lifting for automatic design of path-complete stability certificates. Version arXiv:2607.00637v2 (CC BY 4.0), distributed by arXiv under the Creative Commons Attribution 4.0 International licence, http://creativecommons.org/licenses/by/4.0/, as recorded in the arXiv licence metadata for this exact version and shown on the abstract page https://arxiv.org/abs/2607.00637v2. Full licence notice: \texttt{licenses/arxiv-2607.00637-CC-BY-4.0.txt}.\par\smallskip

\noindent\textit{Editorial scope.} Counted unit: the article's theorem thm:relaxation\_forward (source0.tex lines 420--426). The article's complete original proof of it is delivered with it (source0.tex lines 695--740, one \texttt{proof} environment of 46 lines, which the article places away from the statement). No mathematics is written, repaired or re-derived: the statement and the proof are the article's own lines, in the article's own notation, and no retained line is substituted or re-flowed.  Retained uncounted as the source-local context the proof needs: lines 299--313; lines 502--507.  Nothing was omitted from any retained span, so the package carries no omission record.  No retained line contains a citation, so the wrapper carries no bibliography and no bibliography entry.  No retained line contains a figure environment, an image-inclusion command or drawing code, so no image is delivered and the package declares no asset dependency.  Editorial formatting only, in addition: an AMS \texttt{amsart} preamble that restates the article's own declarations and macros used by the retained lines, namely the environments \texttt{problem}, \texttt{theorem}, \texttt{proposition} on separate counters, together with the standard packages those lines need.  The retained environments print in this document as Problem 1, Theorem 1, Proposition 1 in order of appearance, whereas the article prints the same items with the numbering of its own full document.  The complete native source of the article as distributed in the arXiv e-print for this exact version is bundled unchanged as \texttt{sources/204182/source0.tex}.
\begin{problem}
\phantomsection
Let $G=(S,E)$ be a directed labeled graph on $\Mset$ and let
$\mathcal A=\{A_1,\dots,A_M\}$ be a set of matrices. We consider the
following optimization problem:
\begin{align*}
\begin{aligned}
\inf_{\{P_a\}_{a\in S}, \gamma\ge 0} \quad & \gamma \\[1mm]
\text{s.t.} \quad
& \gamma^2 P_a \succeq A_i^\top P_b A_i, && \forall (a,b,i)\in E, \\[1mm]
& P_a \succ 0, && \forall a\in S .
\end{aligned}
\end{align*}
\label{prob:graph_opt}
\end{problem}

\begin{proposition}
\phantomsection
\label{prop:opt_tight}
Let $G=(S,E)$ be a directed labeled graph, and suppose that \Cref{prob:graph_opt} admits an optimal solution $(\gamma^\star ,\{P_a^\star \}_{a\in S})$ on $G$. Let $\tilde G=(S,\tilde E)$ be the tight subgraph associated with this solution.
Then $\gamma^\star (\tilde G) = \gamma^\star (G)$.
\end{proposition}

\begin{theorem}
\phantomsection
\label{thm:relaxation_forward}
Let $G=(S,E)$ be a directed labeled graph and let $v\in S$.
Let $G^v$ be the forward lift of $G$ with respect to $v$.
Then $\gamma^\star (G^v) \le \gamma^\star (G)$.
\end{theorem}

\begin{proof}
Let $v\in S$ with at most one outgoing edge in $\tilde E$. We already know that $\gamma^\star (G^v)\le\gamma^\star (G)$ by
\Cref{thm:relaxation_forward}. Let
$(\gamma,\{\hat P_a\}_{a\in S^v})$ be feasible for
\Cref{prob:graph_opt} on $G^v$. We construct
$(\gamma,\{P_a\}_{a\in S})$ feasible for $\tilde G=(S,\tilde E)$.

If $\{(v,b,i)\in\tilde E\}=\{(v,\bar b,\bar i)\}$, set
\[
P_a\coloneqq\hat P_a \ (a\neq v),\qquad
P_v\coloneqq\hat P_{v^{(\bar b)}} .
\]
If $\{(v,b,i)\in\tilde E\}=\emptyset$, choose
$\bar b\in\mathrm{Post}(v)$ and set $P_v\coloneqq\hat P_{v^{(\bar b)}}$.

We verify feasibility for each $(a,b,i)\in\tilde E$.

If $a,b\neq v$, the edge appears identically in $G^v$, hence
\[
\gamma^2 P_a\succeq A_i^\top P_b A_i .
\]

If $a=v$ and $b\neq v$, then necessarily $b=\bar b$, and $G^v$
contains the edge $(v^{(\bar b)},\bar b,i)$, so
\[
\gamma^2 P_v=\gamma^2 \hat P_{v^{(\bar b)}}\succeq A_i^\top \hat P_{\bar b} A_i=A_i^\top P_{\bar b} A_i.
\]

If $a=b=v$, then $\bar b=v$ and $G^v$ contains
$(v^{(v)},v^{(v)},i)$, and feasibility of $G^v$ gives
\[
\gamma^2 \hat P_{v^{(v)}}\succeq A_i^\top \hat P_{v^{(v)}} A_i ,
\]
which yields $\gamma^2 P_v\succeq A_i^\top P_v A_i$.

If $a\neq v$ and $b=v$, then $(a,v^{(w)},i)\in E^v$ for any
$w\in\mathrm{Post}(v)$, hence
\[
\gamma^2 \hat P_a\succeq A_i^\top \hat P_{v^{(w)}} A_i .
\]
Taking $w=\bar b$ gives $\gamma^2 P_a\succeq A_i^\top P_v A_i$.

Thus $(\gamma,\{P_a\})$ is feasible for $\tilde G$. Therefore every feasible solution for $G^v$ yields a feasible solution for $\tilde G$ with the same objective value, implying $\gamma^\star (G^v)\ge\gamma^\star (\tilde G)$. Since
$\gamma^\star (G)=\gamma^\star (\tilde G)$ by \Cref{prop:opt_tight},
we obtain $\gamma^\star (G^v)\ge\gamma^\star (G)$, concluding the proof.
\end{proof}

\end{document}
